\subsection{C1 and C3: precise approximation deductions}

Let $A$ be the henselization at a prime of an algebra of finite type over a
field or an excellent DVR, and let $X$ be proper over $A$.
Artin's Theorem~3.5 \cite{Artin1969}, pp.~31--32, says that
\[
 \operatorname{Pic}(X)\longrightarrow\operatorname{Pic}(X_{\widehat A})
\]
is injective with dense image for the topology of restriction to the closed
fibre thickenings. If $X_{\widehat A}$ is projective, approximate an ample
line bundle on it modulo the maximal ideal. The resulting line bundle on
$X$ is ample on the closed fibre. Openness of relative ampleness makes it
relatively ample on a neighbourhood of the closed point of $\operatorname{Spec}A$;
that neighbourhood is the entire local spectrum. Thus $X/A$ is projective.
This proves C1 under the stated approximation hypotheses in all characteristics.

Elkik \cite{Elkik1973}, III~\S1, Theorem~3, p.~569 and its explicit Picard
corollary p.~573, proves the following stronger version of C3. For every
noetherian henselian pair $(A,J)$, with $\widehat A$ its $J$-adic completion,
\[
 \operatorname{Pic}(\operatorname{Spec}A-V(J))
 \xrightarrow{\sim}
 \operatorname{Pic}(\operatorname{Spec}\widehat A-V(J\widehat A)).
\tag{C.1}
\]
Taking $J$ maximal proves C3, without an isolated-singularity or characteristic
zero assumption.

For C2 in characteristic zero, write $\mathfrak m$ for the maximal ideal and
suppose $H^1(X,\mathcal O_X)$ has finite length. The proof of
\cite[Proposition~1.2]{PerlingSchroer2017} shows that
$H^1(X,\mathfrak m^n\mathcal O_X)=0$ for $n\gg0$: the graded direct sum is
finite over the Rees algebra, its sufficiently high graded terms are obtained
by multiplication by $\mathfrak m$, and each such term has finite length.
In characteristic zero formal logarithm identifies
$1+\mathfrak m^n\mathcal O_{\widehat X}$ with
$\mathfrak m^n\mathcal O_{\widehat X}$. Formal functions therefore gives
$H^1(\widehat X,1+\mathfrak m^n\mathcal O_{\widehat X})=0$ for large $n$.
Hence $\operatorname{Pic}(\widehat X)\to\operatorname{Pic}(X_n)$ is injective.
Approximate a given formal line bundle at this one finite level using Artin's
theorem; injectivity then identifies its approximation with the original
formal bundle. This proves C2 in characteristic zero.
Formal logarithm does not apply to an arbitrary maximal ideal in mixed
characteristic. The arguments below replace it first for a DVR and then
in arbitrary dimension. Perling--Schr\"oer's two-cohomology finiteness
theorem will be used only on a suitable auxiliary modification.

\subsection{C2 over a henselian DVR, in every characteristic}

\paragraph{Theorem C2.D.}
Let $R$ be an excellent henselian DVR for which Artin approximation applies,
and let $X$ be any proper $R$-scheme. If $H^1(X,\mathcal O_X)$ has finite
length, then $\operatorname{Pic}(X)\to\operatorname{Pic}(X_{\widehat R})$
is an isomorphism. No flatness of $X$ and no finite-length hypothesis on
$H^2(X,\mathcal O_X)$ are needed.

First suppose $X/R$ flat, and write $\pi$ for a uniformizer. Work after
completion to prove finite determination. Choose $c,d\ge0$ such that
$\pi^c H^1(X,\mathcal O_X)=0$ and $\pi^d$ kills the torsion submodule of
$H^2(X,\mathcal O_X)$; these exist by proper coherent cohomology finiteness.
Set $e=c+d$. The cohomology sequence for multiplication by $\pi^n$ gives
\[
 0\longrightarrow H^1(\mathcal O_X)/\pi^n
 \longrightarrow H^1(\mathcal O_{X_n})
 \longrightarrow H^2(\mathcal O_X)[\pi^n]\longrightarrow0,
 \qquad X_n=V(\pi^n).
\]
Consequently $\pi^e H^1(\mathcal O_{X_n})=0$ for every $n$.
If $L|_{X_n}\simeq\mathcal O_{X_n}$ with $n>e$, the analogous sequence
injects $H^1(L)/\pi^n H^1(L)$ into $H^1(\mathcal O_{X_n})$. Therefore
$\pi^e H^1(L)=\pi^nH^1(L)$, and Nakayama applied to the finite module
$\pi^eH^1(L)$ gives $\pi^eH^1(L)=0$.

Let $a_n=1$ be the trivializing section of $L|_{X_n}$. In the multiplication
exact sequences, the obstruction to lifting its restriction $a_{n-e}$ to
$H^0(X,L)$ is $\pi^e$ times the obstruction for $a_n$; it is zero. Thus
$a_{n-e}$ lifts. Its reduction on the closed fibre is a generator, so the
lift is a nowhere-vanishing section on all of $X$: its failure locus is
closed, and its proper image in the local base would meet the closed point.
Hence $L\simeq\mathcal O_X$. This proves injectivity of
$\operatorname{Pic}(X_{\widehat R})\to\operatorname{Pic}(X_n)$ for $n>e$.
Artin's density theorem now proves surjectivity of completion on Picard
groups, as in C1; its injectivity was already proved there.

For a general proper $X/R$, let $I\subset\mathcal O_X$ be its coherent
$\pi$-power torsion ideal and put $X'=V(I)$. Then $X'$ is flat over $R$.
For $N$ sufficiently large, Artin--Rees and bounded torsion give
$I\cap\pi^N\mathcal O_X=0$. Thus
\[
 \mathcal O_X\simeq
 \mathcal O_{X'}\times_{\mathcal O_{X'_N}}\mathcal O_{X_N}.
\]
The generic fibres of $X$ and $X'$ agree, so $H^1(X',\mathcal O_{X'})$
is a finite torsion $R$-module and the flat case applies to $X'$.
An invertible sheaf on the completed $X$ restricts to one on completed
$X'$, which algebraizes, and one on $X_N$, which is already defined over
$R$. Their gluing isomorphism on $X'_N$ is also defined over $R$, since
$R/\pi^N=\widehat R/\pi^N$. Invertible modules patch across this fibre
product (one of the two maps is surjective), giving a line bundle on $X$.
This completes the proof.

The one-parameter proof uses multiplication by a uniformizer. The following
argument turns an arbitrary maximal ideal into an invertible ideal by a
blowup and supplies the additional descent step.

\subsection{C2 in arbitrary dimension and characteristic}

\paragraph{Theorem C2.G.}
Let $A$ be an excellent henselian local ring, $f:X\to\operatorname{Spec}A$
a proper morphism, and suppose $H^1(X,\mathcal O_X)$ has finite length.
Then
\[
 \operatorname{Pic}(X)\xrightarrow{\sim}
 \operatorname{Pic}(X_{\widehat A}).                 \tag{C2.G.1}
\]
In particular this includes every henselization in Part C, in positive,
zero, and mixed characteristic, without assuming $X/A$ flat.

\paragraph{Proof.}
Artin approximation gives injectivity and density as above; for the stated
henselizations this is Artin's Theorem~3.5, and for general excellent
henselian rings the same argument uses Popescu approximation, as explained
in \cite[Remark~1.3]{PerlingSchroer2017}. We prove the required finite
determination after completion. We use formal functions in the precise
form of Stacks tag~02OB(3): for a proper $A$-scheme $W$ and every fixed
$k$, there is $N_0$ such that for $N\geq N_0$,
\[
 \operatorname{Im}\bigl(H^0(W_N,\mathcal O)\to H^0(W_k,\mathcal O)\bigr)
 =\operatorname{Im}\bigl(H^0(W,\mathcal O)\to H^0(W_k,\mathcal O)\bigr),
 \quad W_n=V(\mathfrak m^n\mathcal O_W).             \tag{C2.G.2}
\]
An element on the right lifting a unit on the closed fibre is a global unit:
its nonunit locus is closed with proper image in the local base, hence
would meet the closed fibre if it were nonempty.

First remove the coherent ideal $I\subset\mathcal O_X$ of sections killed
by a power of $\mathfrak m$. A uniform power kills $I$, and Artin--Rees
gives $I\cap\mathfrak m^N\mathcal O_X=0$ for $N$ sufficiently large.
With $X^\circ=V(I)$ we have the fibre-product algebra
\[
 \mathcal O_X=
 \mathcal O_{X^\circ}\times_{\mathcal O_{X^\circ_N}}\mathcal O_{X_N}.
\tag{C2.G.3}
\]
The cohomology sequence and the fact that all cohomology of $I$ is killed
by a power of $\mathfrak m$ show that $H^1(X^\circ,\mathcal O)$ still has
finite length. If (C2.G.1) holds for $X^\circ$, an arbitrary completed line
bundle patches from its algebraized restriction to $X^\circ$ and its
restriction to $X_N$. Their gluing isomorphism is defined on
$X^\circ_N$, which is unchanged by completion. Invertible modules patch
over a fibre-product ring with a surjective side map (Stacks
Lemma~15.6.9, also used in tag~0DVP). It suffices therefore to handle $X$
with no $\mathfrak m$-power torsion in its structure sheaf.

From now on $A$ is complete and $X$ has this property. Put
\[
 Z=\operatorname{Bl}_{\mathfrak m}(\operatorname{Spec}A),\qquad
 X'=\operatorname{Bl}_{\mathfrak m\mathcal O_X}(X),\qquad
 \pi:X'\to X.
\]
The blowup makes $\mathfrak m\mathcal O_{X'}$ an invertible ideal defining
an effective Cartier divisor (Stacks tag~02OS), and its universal property
(tag~0806) gives a proper map $X'\to Z$. Take its Stein factorization
\[
 X'\xrightarrow{h}T\longrightarrow Z.
\]
Here $T\to Z$ is finite, $T/A$ is proper, and
$h_*\mathcal O_{X'}=\mathcal O_T$ (tag~03H0). The ideal
$J=\mathfrak m\mathcal O_T$ is invertible and locally generated by a
nonzerodivisor: a generator acts injectively on $\mathcal O_{X'}$, and
therefore on its pushforward.

Away from $V(J)$, the map $X'\to Z$ is the original $X$ over the punctured
base. Write $g:T\to Z$ for the finite morphism. Since $g_*$ is exact and
faithful on quasi-coherent modules, the identity
$g_*R^1h_*\mathcal O_{X'}=R^1(g\circ h)_*\mathcal O_{X'}$ and the
finite-length hypothesis imply that $R^1h_*\mathcal O_{X'}$ is supported
on $V(J)$, and hence is killed
by $J^c$ for some $c$. The $J$-power torsion in
$R^2h_*\mathcal O_{X'}$ is killed by $J^d$ for some $d$.
Set $e=c+d$.

The argument of Theorem C2.D applies on every affine chart of $T$ where
$J=(a)$, after $a$-adic completion. It only requires that $a$ be a
nonzerodivisor on $X'$ and that $a^cH^1(\mathcal O)=0$; the base need not
be a DVR. The use of Nakayama is valid because $a$ belongs to the Jacobson
radical of its adic completion. Thus, if a line bundle $L$ on $X$ has a
chosen trivialization on $X_N$ with $N>e$, its pullback has, locally on
the formal completion of $T$, generators agreeing with that trivialization
modulo $J^{N-e}$.

The ratios of these generators descend to units on formal $T$: formal
functions and $h_*\mathcal O_{X'}=\mathcal O_T$ identify the completed
pushforward of the structure sheaf with $\mathcal O_{\widehat T}$.
They define a formal line bundle $M$ on $T$ with
$h^*M=(\pi^*L)^{\wedge}$. Moreover $M$ is trivial modulo $J^{N-e}$.
Indeed the local multiplication sequence gives an injection
\[
 \mathcal O_T/J^r\longrightarrow h_*\mathcal O_{X'_r};
\]
so a descended ratio equal to one on $X'_r$ is already one on $T_r$.

Off the closed point of $\operatorname{Spec}A$, the scheme $T$ is the
finite Stein space $\operatorname{Spec}(f_*\mathcal O_X)$. Hence
$H^1(T,\mathcal O_T)$ and $H^2(T,\mathcal O_T)$ both have finite length.
Formal GAGA first algebraizes $M$ on the proper $A$-scheme $T$.
Apply \cite[Proposition~1.2 and the proof of Theorem~1.1]{PerlingSchroer2017}
with the structure sheaf: for some $q$, a formal line bundle on $T$ trivial
on $T_q$ is trivial. Choose $N\geq q+e$. Then $M$ is trivial, and formal
GAGA over complete $A$ implies $\pi^*L\simeq\mathcal O_{X'}$.
Notice that finite length of $H^2$ has been used only on this auxiliary
$T$, where it follows from the construction, not on the original $X$.

It remains to descend triviality at finite order. Let
$Y=\operatorname{Spec}_X(\pi_*\mathcal O_{X'})$ be the finite Stein space
of $\pi$. The inclusion $\mathcal O_X\hookrightarrow\mathcal O_Y$ is
injective: its kernel is zero off the closed fibre, and $X$ has no
$\mathfrak m$-power torsion. Also $Y\to X$ is an isomorphism off that
fibre. The map $\operatorname{Pic}(Y)\to\operatorname{Pic}(X')$ is
injective by the projection formula and the defining Stein equality.
Therefore $L|_Y$ is trivial.

Let $C\subset X$ and $C'\subset Y$ be the conductor subschemes. Their
ideals contain $\mathfrak m^k$ for some $k$, and
$\mathcal O_X=\mathcal O_Y\times_{\mathcal O_{C'}}\mathcal O_C$.
A trivialization of $L|_Y$ and the chosen trivialization of $L|_{X_N}$
differ on $Y_N$ by a unit $u_N$. Increase $N$ to satisfy (C2.G.2) for
$W=Y$ and this $k$. Then $u_N|_{Y_k}$ lifts to a global function on $Y$,
which is a unit by the observation following (C2.G.2). Adjusting the
trivialization on $Y$ by that unit makes it agree with the one on $C$
after pullback to $C'$. Fibre-product patching gives $L\simeq\mathcal O_X$.

All bounds depend only on $X/A$, not on $L$. This proves finite
determination in the case with no closed-fibre torsion. Artin's density
then gives algebraization in that case, and (C2.G.3) gives it in general.
Together with injectivity, this proves (C2.G.1).

\subsection{The characteristic-$p$ local NS comparison: a discrete obstruction}

The original setting of 57:184--194 is the local ring at the vertex of an
affine cone over a smooth projective variety $V$, with its given very ample
bundle $\mathcal O_V(1)$. In that setting the characteristic-$p$ assertion
has a direct proof in every dimension. The blowup is the total space of
$\mathcal O_V(-1)$, and its completion along the zero section has a projection
back to $V$. Thus restriction to $V$ on Picard groups is split surjective.
The kernel has the filtration arising from $1+I^n$, with successive quotients
subquotients of $H^1(V,\mathcal O_V(n))$. Both $H^1$ and $H^2$ of these sheaves
vanish for $n\gg0$, so the formal Picard group is determined at a finite level.
In characteristic $p$ these finitely many successive additive groups are killed
by $p$. The comparison cokernel is consequently killed by some $p^N$.
Removing the exceptional divisor quotients both groups by the same class
$\mathcal O_V(1)$, and the conclusion persists. Elkik then identifies the
henselian and completed punctured Picard groups. This proves the cone-context
claim even at the level of Picard groups modulo bounded $p$-power torsion;
it is not a proof of the extrapolation to all local schemes.

Under the literal comparison of the Picard groups of the three punctured
spectra, H12 is false, already in characteristic three and dimension three.
Let $k$ be algebraically closed of characteristic three and put
\[
 A=\left(k[x,y,z,t]/(xy-z^2-t^2(1+t))\right)_{(x,y,z,t)},
 \qquad U=\operatorname{Spec}A-\{\mathfrak m\}.
\tag{C.2}
\]
This is a normal isolated hypersurface singularity of dimension three.
The polynomial $z^2+t^2(1+t)$ is irreducible in $k[z,t]$: a factorization
would make $1+t$ a square in $k(t)$, impossible because it has an odd valuation
at $t=-1$. Thus $(x)$ is a prime divisor in $A$. The ring $A[1/x]$ is a
localization of $k[x,x^{-1},z,t]$, hence factorial. Nagata's class-group
localization sequence then says that $\operatorname{Cl}(A)$ is generated by
the principal prime $(x)$, so $\operatorname{Cl}(A)=0$. Since $U$ is regular,
$\operatorname{Pic}(U)=0$.

Both the henselization and the completion contain the unique square root
$s=\sqrt{1+t}$ congruent to one. For $i^2=-1$, set $u=z-its$ and $v=z+its$.
Their linear terms are independent because $2i\ne0$. The formal inverse
function theorem gives
\[
 \widehat A\simeq k[[x,y,u,v]]/(xy-uv).
\]
The divisor class group of this completed threefold node is $\mathbf Z$,
generated by $(x,u)$. One direct computation is to blow up its closed point:
the exceptional divisor is $Q=\mathbf P^1\times\mathbf P^1$, its normal bundle
is $\mathcal O_Q(-1,-1)$, and the successive conormal powers are
$\mathcal O_Q(n,n)$. Their $H^1$ and $H^2$ vanish for $n\geq1$, so formal
functions and existence identify the Picard group of the blowup with
$\operatorname{Pic}(Q)=\mathbf Z^2$. Removing the exceptional divisor quotients
by its class $(-1,-1)$, leaving $\mathbf Z$. Formula (C.1) gives the same
punctured Picard group for $A^h$.

There is no connected algebraic part accounting for this extra integer:
the node is Cohen--Macaulay of dimension three, so the tangent space of its
local Picard functor is $H^1(U,\mathcal O_U)=H^2_{\mathfrak m}(A)=0$.
Equivalently the above exceptional-quadric computation has zero connected
Picard part. Thus the comparison is $0\to\mathbf Z\to\mathbf Z$ on the
resulting discrete local divisor groups, and its cokernel is not a $3$-group.
Some modern definitions define the local Picard/NS group using a strict
henselization from the outset; with that convention the notation
$\mathrm{NS}_{\mathrm{loc}}(S)$ already incorporates the second ring, and
one must not interpret it as $\operatorname{Pic}(U)/\operatorname{Pic}^0(U)$.
The ambiguity must be resolved before identifying this with a differently
defined local N\'eron--Severi invariant.

